\section{A smooth counterexample to integration-grid agreement}
\label{sec:aliasing}

\subsection{Unresolved coordinate variation}
\label{sec:aliasing-proof}

For an integer $K>0$, define
\begin{align}
h(\alpha)&=\frac25\sin^2(\pi K\alpha),\\
g(x_1,x_2)&=\frac25x_1+\frac35x_2+
(x_1-x_2)h\!\left(\frac{x_1+x_2}{2}\right),\\
Q(x_1,x_2)&=-\frac12\bigl(g(x_1,x_2)-1\bigr)^2.
\end{align}
Take baseline $(0,0)$, actual input $(1,1)$, and fixed reference $g(1,1)=1$.
The transverse term vanishes along $(\alpha,\alpha)$, so
$g(\alpha,\alpha)=\alpha$ and $Q(\alpha,\alpha)=-(1-\alpha)^2/2$.
Differentiating before restricting to that path gives
\begin{equation}
\nabla Q(\alpha,\alpha)=(1-\alpha)
\left(\frac25+h(\alpha),\frac35-h(\alpha)\right).
\label{eq:aliasing-gradient}
\end{equation}
Thus the oscillation transfers gradient between coordinates while leaving their
sum unchanged.

Use the composite trapezoidal rule with $m$ intervals.
Whenever $m$ divides $K$, every node $\alpha=j/m$ has $h(\alpha)=0$.
The sampled coordinate functions are then linear. Both the trapezoidal rule and endpoint-inclusive average integrate a linear function exactly on any uniform grid, giving
\begin{equation}
A_m^{\mathrm{trap}}=\left(\frac15,\frac3{10}\right),\qquad
\mathbf{1}^{\mathsf T}A_m^{\mathrm{trap}}=\frac12=Q(1,1)-Q(0,0).
\label{eq:aliasing-coarse}
\end{equation}
Here $\mathbf{1}=(1,1)^{\mathsf T}$ sums the coordinates. Every such grid has zero
completeness residual and ranks feature 2 first, giving identical masks and
perturbation curves under a deterministic ranking rule.

The exact coordinate integrals differ. Since $h(1-\alpha)=h(\alpha)$,
substitution gives
\[
\int_0^1\alpha h(\alpha)\,d\alpha
=\int_0^1(1-\alpha)h(\alpha)\,d\alpha.
\]
Adding these expressions and using $\int_0^1 h(\alpha)\,d\alpha=1/5$
shows that either integral is $1/10$. Consequently,
\begin{equation}
A_{\mathrm{exact}}^Q=\left(\frac3{10},\frac15\right).
\label{eq:aliasing-exact}
\end{equation}
The exact order puts feature 1 first. Coarse errors $(-1/10,1/10)$ cancel in
completeness, but their $L_1$ norm, $1/5$, is 40\% of the exact attribution's norm.

On the $2K$ grid, $h$ alternates between zero and $2/5$, with trapezoidal integral
$1/5$. Reflection symmetry gives area $1/10$ for $(1-\alpha)h$, recovering
Equation~\ref{eq:aliasing-exact}. This recovery is construction-specific. For
finitely many uniform budgets, a common multiple $K$ hides every oscillation.
Agreement cannot certify coordinate integrals without a regularity bound
controlling unresolved variation.

The endpoint-inclusive average has the same coarse failure.
When $m$ divides $K$, the average of the sampled linear functions also yields
Equation~\ref{eq:aliasing-coarse}. On the $2K$ grid, only odd nodes contribute
to $(1-\alpha)h$. Their sum is
$\frac25\sum_{\ell=0}^{K-1}(1-(2\ell+1)/(2K))=K/5$.
Dividing by the $2K+1$ nodes gives
\begin{equation}
A_{2K}^{\mathrm{avg}}=
\left(\frac15+\frac{K}{5(2K+1)},\
\frac3{10}-\frac{K}{5(2K+1)}\right).
\label{eq:aliasing-average}
\end{equation}
The first coordinate exceeds the second for $K>1/2$, including every allowed
integer. Coordinates remain inexact; exact $2K$ recovery belongs to the trapezoidal rule.

\subsection{Exact IG and perturbation quality}
\label{sec:aliasing-quality}

Set $K=128$, so $h(1/2)=0$ and the whole-feature interventions give
$g(1,0)=2/5$ and $g(0,1)=3/5$. Under the same normalized Q response
and fractions $\{0,1/2,1\}$, feature-2-first insertion values are
$(0,0.84,1)$ and deletion values are $(1,0.64,0)$.
Their trapezoidal AUCs are 0.67 and 0.57.
The exact-IG feature-1-first order instead gives insertion AUC 0.57 and deletion AUC 0.67.
The inaccurate order therefore scores better by both conventional criteria.
The transverse term alters coordinate gradients along the diagonal without
altering these finite off-path intervention values. Exact path integration
and perturbation-score optimization are distinct questions.

CPU fp32/fp64 checks at $K=128$ verify trapezoidal coordinates, completeness,
repeatability and responses: budgets 32, 64 and 128 alias; 256 recovers the analytic
order to numerical tolerance. The endpoint-inclusive result is analytic only;
these checks use no RDT data.
